% Exact source-aware P06.2 candidate packet, 25 August 2026.
% No novelty or priority claim is intended.

\paragraph{Source scope and conventions.}
We use the CLD current convention: on a pure $n$-dimensional analytic
space, a current of bidegree $(1,1)$ acts on smooth compactly supported
forms of bidegree $(n-1,n-1)$.  All current identities below are local and
all test forms have compact support in the relative interior of the compact
analytic domain on which a tropicalization is used.

The primary sources used here are the following.
\begin{enumerate}
\item Gubler--Jell--Rabinoff, \emph{Forms on Berkovich spaces based on
harmonic tropicalizations}, arXiv:2111.05741v3.  Proposition~5.8 writes an
$\mathbf R$-PL function, after shrinking in the analytic topology, as a
finite real linear combination of integral PL functions.  Lemma~5.10 lets
one choose the coefficients $\mathbf Q$-linearly independent.
Proposition~8.3 and paragraph~8.4 show that a finite tuple of integral PL
functions is locally covered by one smooth tropicalization, with
piecewise $(\mathbf Z,\Gamma)$-affine coordinate functions.  The local
calculation in the proof of Proposition~9.6 gives, after adding an auxiliary
moment map which separates the chosen compact support from the analytic
boundary, the identity between the analytic curvature pairing and the
tropical divisor pairing.  The proof explicitly uses CLD's tropical
Poincare--Lelong formula.  Proposition~7.1 and Definition~7.2 identify
harmonicity of an $\mathbf R$-PL germ with numerical triviality of its
residue class.
\item Chambert--Loir--Ducros, \emph{Formes diff\'erentielles r\'eelles et
courants sur les espaces de Berkovich}, arXiv:1204.6277v2, 25 July 2025.
The construction of smooth forms and currents by tropical charts, together
with the sheaf property and smooth partitions of unity, permits a compactly
supported test form to be split into finitely many forms tropicalized on
compact domains.  The tropical integration formula, Green's formula, and
tropical Poincare--Lelong identify the curvature of an integral PL function
with the integration current of its tropical Weil divisor in such a chart.
Theorem~18.8.4, with $d=1$, says for an integral PL metric on a
boundaryless separated analytic space that its curvature current is zero
near a point if and only if the first Chern class of its residual line
bundle is numerically trivial there.
\item Gubler--Rabinoff, \emph{Classically psh and pluriharmonic functions
on Berkovich spaces}, arXiv:2506.13548v2.  Theorem~4.3 says, over a
nontrivially valued field, that an $\mathbf R$-PL function on a good strict
analytic space is classically psh if and only if it is semipositive in the
residue-nef sense.
\end{enumerate}
For source identity, the audited PDF SHA-256 values are respectively
\texttt{572d51c9f8907dbb814684f7d136612243c0e0b70cddbe6759bb686f0bcfbd42},
\texttt{5690f999523017d184ab751bb55857ab1637b197d1186ad4fd2edba7e1bda719},
and
\texttt{83ca9a19394f44bb2e4e1aa316de95938f311acd53a00fe250fa0f605184b726}.

\begin{proposition}
Let $K$ be a complete nontrivially valued non-Archimedean field, and let
$X$ be a good, separated, boundaryless, strictly $K$-analytic space that is
pure of dimension $n\geq1$.  If $h:X\to\mathbf R$ is an $\mathbf R$-PL function and
its Chambert--Loir--Ducros current satisfies
\[
                         \mathrm d'\mathrm d''[h]=0,
\]
then $h$ is harmonic in the sense of Gubler--Jell--Rabinoff.  Equivalently,
the residue class $L_h(x)\in\operatorname{Pic}(\operatorname{red}(X,x))_
{\mathbf R}$ is numerically trivial for every $x\in X$.  Consequently $h$
is classically pluriharmonic in the sense of Gubler--Rabinoff.
\end{proposition}

\begin{proof}
Fix $x\in X$.  By Gubler--Jell--Rabinoff, Proposition~5.8 and
Lemma~5.10, after replacing $X$ by an analytic neighbourhood $W$ of $x$,
there are integral PL functions $h_1,\ldots,h_s$ on $W$ and real numbers
$\lambda_1,\ldots,\lambda_s$ which are linearly independent over
$\mathbf Q$ such that
\begin{equation}
\label{eq:p06-independent-decomposition}
                         h=\sum_{i=1}^{s}\lambda_i h_i
                         \qquad\text{on }W.
\end{equation}
We first prove
\begin{equation}
\label{eq:p06-integral-current-vanishing}
                         \mathrm d'\mathrm d''[h_i]=0
                         \qquad(1\leq i\leq s)
\end{equation}
as currents on a possibly smaller neighbourhood of $x$.

Here is the coefficient-separation argument, including the refinement
needed to prevent cancellation in a coarse tropical chart.  Let $n$ be the
pure dimension near $x$, and let $\alpha$ be an arbitrary smooth compactly
supported $(n-1,n-1)$-test form in a sufficiently small open subset of
$W$.  By the sheaf property for currents and a smooth partition of unity,
it is enough to treat one summand for which there are a compact strictly
analytic domain $U\subset W$, a smooth tropicalization
\[
                         H:U\longrightarrow\mathbf R^q,
\]
and a smooth Lagerberg form $\beta$ on $H(U)$ such that
$\alpha=H^*\beta$ and $\operatorname{Supp}(\alpha)\subset\operatorname{Int}(U)$.
Shrinking the summand once more if necessary, the boundary-separation
construction used in the proof of Gubler--Jell--Rabinoff,
Proposition~9.6, supplies an auxiliary smooth tropicalization $G_0$ and a
smooth cutoff on the target so that, after multiplying $\beta$ by that
cutoff, its support lies in
\[
       \Omega:=G(U)\setminus G(\partial U),
\]
without changing its pullback on the support of $\alpha$.  Here $G$ is the
common refinement constructed below.  This is precisely the
compact-support/boundary-avoidance step in that published local
calculation; it ensures that every pulled-back target form used below is a
legitimate compactly supported analytic test form.

Apply Proposition~8.3 and paragraph~8.4 of the same paper simultaneously
to the tuple $(h_1,\ldots,h_s)$ on $U$.  It gives a smooth tropicalization
$F:U\to\mathbf R^m$ and piecewise $(\mathbf Z,\Gamma)$-affine functions
$\varphi_i$ on $F(U)$ with $h_i=\varphi_i\circ F$.  Replace all charts by
the common smooth tropicalization
\[
               G=(F,H,G_0):U\longrightarrow\mathbf R^N
\]
and put $\psi_i=\varphi_i\circ\operatorname{pr}_F$ on $G(U)$.  Thus
$h_i=\psi_i\circ G$, while the chosen test form is $G^*\omega$ for a smooth
Lagerberg $(n-1,n-1)$-form $\omega$ compactly supported in $\Omega$.

Endow the $n$-dimensional part of $G(U)$ with its canonical tropical
weights.  These weights are integers, and each $\psi_i$ has integral
slopes.  The local current calculation in the proof of
Gubler--Jell--Rabinoff, Proposition~9.6, applied now to the assumed current
identity rather than to a previously known harmonic function, gives for
every smooth target test form $\eta$ compactly supported in $\Omega$
\begin{align}
  0
   &=\bigl\langle\mathrm d'\mathrm d''[h],G^*\eta\bigr\rangle\notag\\
   &=\left\langle
       \mathrm d'\mathrm d''
       \left[\sum_{i=1}^{s}\lambda_i\psi_i\right],\eta
     \right\rangle_{G_*(U)} .
\label{eq:p06-push-to-common-chart}
\end{align}
No pullback of the current $\mathrm d'\mathrm d''[h]$ is being taken in
\eqref{eq:p06-push-to-common-chart}: the first pairing is simply evaluation
against the smooth analytic test form $G^*\eta$, and the second equality is
the defining tropical integration formula for that pairing.

The tropical Poincare--Lelong formula identifies the target current in
\eqref{eq:p06-push-to-common-chart} with the integration current of
\[
       \operatorname{div}\left(\sum_i\lambda_i\psi_i\right)
       =\sum_i\lambda_i\operatorname{div}(\psi_i)
\]
on the weighted tropical cycle $G_*(U)$ over $\Omega$.  After a common
polyhedral subdivision, the coefficient of every codimension-one face in
$\operatorname{div}(\psi_i)$ is an integer: it is a sum of products of
integer tropical multiplicities and integral normal slopes.  Vanishing of
the integration current forces the coefficient of each such face to
vanish.  Hence, face by face,
\[
                         \sum_i\lambda_i a_i=0,
                         \qquad a_i\in\mathbf Z.
\]
The $\mathbf Q$-linear independence of the $\lambda_i$ implies
$a_i=0$ for all $i$.  Therefore
\begin{equation}
\label{eq:p06-each-target-divisor-zero}
                         \operatorname{div}(\psi_i)=0
                         \quad\text{over }\Omega
                         \qquad(1\leq i\leq s).
\end{equation}

Pairing \eqref{eq:p06-each-target-divisor-zero} with the particular target
form $\omega$ representing $\alpha$ and applying the same tropical
integration formula gives
\[
       \bigl\langle\mathrm d'\mathrm d''[h_i],\alpha\bigr\rangle
       =\bigl\langle[\operatorname{div}(\psi_i)],\omega\bigr\rangle
       =0.
\]
The crucial point is that $G$ contains $H$, the chart tropicalizing this
arbitrary $\alpha$.  Thus no equality obtained after a fixed coarse
pushforward is being pulled back, and possible cancellation between fibres
identified by a previous tropicalization cannot affect the conclusion.
Since the localized summands span the smooth compactly supported test
forms, this proves \eqref{eq:p06-integral-current-vanishing}.

For each $i$, apply Chambert--Loir--Ducros, Theorem~18.8.4 with $d=1$ to
the trivial line bundle carrying the integral PL metric determined by
$h_i$.  The space is separated and boundaryless, exactly as required by
that theorem.  In view of
\eqref{eq:p06-integral-current-vanishing}, it follows that
$L_{h_i}(x)$ is numerically trivial.  Residues depend linearly on PL
functions, so \eqref{eq:p06-independent-decomposition} yields
\[
                         L_h(x)=\sum_i\lambda_iL_{h_i}(x),
\]
which is numerically trivial.  Since $x$ was arbitrary,
Gubler--Jell--Rabinoff, Proposition~7.1 and Definition~7.2 show that $h$
is harmonic.

Finally, harmonicity means that both $h$ and $-h$ are semipositive.  By
Gubler--Rabinoff, Theorem~4.3, both are classically psh.  Thus $h$ is
classically pluriharmonic.
\end{proof}

\paragraph{Scope firewall.}
The proof uses integral-slope arithmetic after decomposing an
$\mathbf R$-PL function.  It gives no assertion for an arbitrary continuous
function.  It also never restricts or pulls back an arbitrary CLD current
to a curve.  Boundarylessness, separation, strictness, goodness, and the
compact-support refinements are retained exactly where used.
